\begin{textsample}{NEG Dim 3 – pi – Score -3.00 – pi/pi\_em\_36.txt}  \label{ex:f3_neg_025}
Língua Inglesa O ensino de Língua Inglesa no Brasil teve início em meados do século XIX, ocupando, a partir de então, a posição de idioma clássico de estudo obrigatório, idioma este que ao longo dos anos sofreu inúmeras mudanças no que diz respeito ao espaço e obrigatoriedade de oferta nas grades curriculares das escolas regulares.
Em 1996, quando a nova LDB/96 foi promulgada, o texto do Art. 26, § 5º passa a dispor que na parte diversificada do currículo será incluído obrigatoriamente o ensino de, pelo menos uma língua estrangeira moderna, ficando a escolha a cargo da comunidade escolar, alinhado às possibilidades das redes e instituições de Ensino Médio de caráter obrigatório, permitindo aos sistemas de ensino, em caráter optativo ofertar outras línguas estrangeiras, preferencialmente o \textbf{Espanhol}.
Assim, a Língua Inglesa apresenta -se, desde a LDB com a reforma do Ensino \textbf{médio}, componente de oferta obrigatória aos currículos dos sistemas de ensino.
A oferta do idioma como componente curricular obrigatório no Ensino Médio possibilita o atendimento das atuais necessidades de colocar o estudante da educação básica em contato com uma língua, que politicamente é considerada língua global, de comunicação internacional cuja aprendizagem é instrumento para aquisição de conhecimentos, exercício da cidadania e símbolo de participação em um mundo plurilíngue e globalizado.
O mapa de habilidades e objetivos propostos para o ensino do componente pressupõe o desenvolvimento de habilidades que ampliam a consciência crítica e reflexiva do estudante, além de prepará -lo para atuar nos diversos campos, sejam culturais, de estudo ou pesquisas, de modo a aprofundar sua compreensão sobre o mundo.
Isso, numa visão intercultural e desterritorializada, considerando as práticas sociais no mundo digital.
O componente curricular de Língua Inglesa foi definido pela Base Nacional Comum Curricular ( BNCC – homologada em 2017 ) como o idioma estrangeiro obrigatório a ser ensinado para a Educação Básica a partir do Ensino Fundamental II de todas as escolas brasileiras.
Assim, apresenta as habilidades específicas do componente considerando que os alunos do Ensino Fundamental devem dominar e que são aprofundadas na última etapa da Educação básica.
A Língua Inglesa é Integrante da área de Linguagens, que tem em sua composição 07 \textbf{competências} e 27 habilidades no Ensino Médio, como se destaca na \textbf{competência} a seguir : ( 1 ) Compreendida como língua de uso mundial, pela multiplicidade e variedade de usos, usuários e funções, o que a caracteriza como língua franca, uma vez que através dela, embora as pessoas não dominem, elas conseguem se comunicar e se fazerem ser compreendidos pelo receptor na contemporaneidade ( BNCC, 2018, p. 476 ).
Isso torna o inglês um componente legitimado pela Base como língua franca uma oportunidade de acesso ao mundo globalizado, deixando de ser apenas dos falantes nativos, o que favorece o ensino da língua com mais interculturalidade, porque varia com diferentes contextos.
O texto da Base define ainda como objetivo para o ensino de língua Inglesa que os estudantes devem desenvolver \textbf{competências} e habilidades a partir de uma perspectiva de educação linguística consciente, crítica e reflexiva, de forma que a aprendizagem do idioma propicie aos estudantes o acesso a novos percursos de construção de conhecimento e o exercício da cidadania ativa, permitindo -lhes vivenciar “ novas formas de engajamento e participação em um mundo social cada vez mais globalizado e plural ” ( BNCC, 2018, p. 239 ).
Assim, o ensino de Língua Inglesa está ancorado na perspectiva de uma abordagem voltada para a oralidade e para o seu uso em um viés que transcende uma aprendizagem voltada para estruturas linguísticas, passando ao universo das estruturas sociais e de aspectos culturais : “ ( 2 ) De acordo com a BNCC considerar que há \textbf{juventudes} implica organizar uma escola que acolha as diversidades e que reconheça os jovens como seus interlocutores legítimos sobre currículo, ensino e aprendizagem ” ( BNCC, 2017, p. 463 ).
Dessa forma, esse texto que apresenta a contextualização para o ensino de Língua Inglesa no Ensino Médio está sequenciado do Contexto Histórico da oferta e os Marcos legais que fundamentam o ensino da língua no país e no estado do Piauí, apresentando uma abordagem crítica da inserção da língua no currículo, sendo considerada inicialmente como língua estrangeira moderna e, posteriormente assumindo o status de língua franca, pois permite a integração com grupos multilíngues e multiculturais no mundo global.
Em seguida são apresentados aspectos que constituem a aprendizagem da língua inglesa e a organização curricular com abordagem das \textbf{competências} e habilidades gerais e específicas da língua para o Ensino Médio.
Assim, espera -se que estudantes, professores e a comunidade escolar piauiense façam uso da língua em diferentes contextos interacionais e comunicativos, possibilitando o compartilhamento de experiências e abrindo caminhos para que o aprendizado e a aplicação da língua aconteçam de forma natural e fruitiva e ainda, que seja meio para ampliar o engajamento discursivo nas redes sociais e demais softwares que permitam as trocas de informação no mundo interativo e globalizado.

% matched lemmas: competência, espanhol, juventude, médio
\end{textsample}
